% Abstract for NavigatorCrew Paper
% XeLaTeX + bidi + IEEE format

\selectlanguage{hebrew}


מאמר זה מציג מערכת ניווט ביומימטית לרחפן רובוטי רך תת-ימי, בהשראת מערכת העצבים המבוזרת של התמנון. המערכת משלבת ארכיטקטורת \en{SLAM} מבוזרת המבוססת על גרף פקטורים, מיזוג חיישנים רב-מודאלי הכולל סונאר, ראייה סטריאוסקופית, חישת מישוש קיבולית ו-\en{IMU}, ותכנון תנועה אדפטיבי באמצעות למידת חיזוק מבוזרת מרובת-סוכנים. המודל הקינמטי מבוסס על קינמטיקת \en{Piecewise Constant Curvature (PCC)} לרובוט רציף, עם דינמיקה מבוססת אלמנטים סופיים וכוחות הידרודינמיים. ארכיטקטורת ה-\en{SLAM} המבוזרת, בהשראת הגנגליונים הפריפריאליים בזרועות התמנון, מאפשרת לכל זרוע לבצע לוקליזציה ומיפוי מקומיים תוך תקשורת מבוזרת באמצעות אלגוריתם \en{ADMM} לאיחוי קונצנזוס. מיזוג החיישנים הרב-מודאלי מבוסס על פילטר חלקיקים מותאם עם שקלול מבוסס אי-ודאות, המאפשר המשך ניווט מדויק גם בנשירת חיישנים. המערכת הוערכה בשלושה תרחישי סימולציה מייצגים: מערה תת-ימית, שדה אצות ומבנה תעשייתי. התוצאות מראות שיפור ממוצע של 36\% ב-\en{ATE (Absolute Trajectory Error)} לעומת \en{SLAM} ריכוזי, עם שיעור הצלחה ממוצע של 93\% וזמן חישוב של 19 אלפיות-שנייה לפריים על פלטפורמת \en{ARM} משובצת. השיטה המוצעת מדגימה את הפוטנציאל של השראה ביולוגית ממערכת העצבים המבוזרת של התמנון להתמודדות עם אתגרי ניווט ייחודיים של רובוטים רכים בסביבות תת-ימיות מוגבלות.


\selectlanguage{hebrew}